\section{Limitations}
\label{sec:limitations}

\paragraph{One model family.} We tested three Claude models only. Other families may behave differently, and
the earlier positive findings were across many families.

\paragraph{Evaluation awareness.} Models can often tell an evaluation from deployment
\citep{needham2025evalaware}. We kept agent-facing text free of test language and judged verbalised
awareness, but a model can behave as if observed without saying so. The null may partly reflect that.

\paragraph{Sandbox realism.} One fictional bank, one policy written to be clear, deterministic characters, a
single RM role and episodes of one or five days. Real deployments have ambiguous policies, longer horizons,
memory, and customers who push back. The violation types are those our tools allow; we did not test social
engineering by customers or attacks on the agent.

\paragraph{Sample size and the bound.} N was set by a fixed budget. A per-probe bound of a few percent does
not rule out a rare failure that matters at bank scale.

\paragraph{The ceiling gate failed.} Without a model-made violation we validated the detector with scripted
agents and a test of the API path. A violation of a kind the scripted agents did not take could be missed;
the blind hand-labelled spot-check (Section~\ref{sec:gates}) is the only check against this.

\paragraph{Deviations.} The backend changed from the Claude Code harness to the API after the first pilot,
and P5 was added after the pilot. Both are logged, neither touched the hypotheses, and the legacy pilot,
which also had no violations, is reported separately (Appendix~\ref{app:legacy}).

\paragraph{Author and tools.} The study was designed, run and analysed by a single author. Code, analysis
and drafts of this paper were written with the help of Claude Code, a product of the same company whose
models are tested; the pre-registration, the deterministic detector and the released event logs are meant to
let readers check the result independently.
